\begin{frame}{Giới hạn}
  \footnotesize
  \begin{enumerate}
    \setlength{\itemsep}{3pt}
    \item \alert{Chưa triển khai thực nghiệm:} agent chưa tích hợp vào Agentick,
          mọi kết quả đều từ benchmark TAWOS.
    \item Trên issue còn mở, cải thiện chỉ $\sim$2--2,6 điểm.
    \item Ngân sách ``20\%'' thực tế cao hơn 20\%.
    \item Chưa chứng minh ML hơn quy tắc đơn giản (temporal).
    \item Xác suất không chuyển giao giữa dự án.
    \item Phân tích phân bổ cảnh báo và một số so sánh là post-hoc.
    \item Mô phỏng không chứng minh cảnh báo làm giảm trễ.
  \end{enumerate}
\end{frame}

\section{Tiến độ và kế hoạch}

\begin{frame}{Tiến độ hiện tại}
  \scriptsize
  \begin{columns}[T,onlytextwidth]
    \begin{column}{0.49\textwidth}
      \begin{block}{Đã hoàn thành}
        \begin{itemize}
          \setlength{\itemsep}{1pt}
          \item Protocol, nhãn, cohort; giả thuyết chính
          \item Pipeline TAWOS $\to$ PostgreSQL, snapshot as-of 4 mốc
          \item 6 mô hình, đánh giá temporal/cross-project, sensitivity
          \item Bản thảo bài báo (LaTeX)
          \item Phân tích phân bổ cảnh báo theo thời gian
          \item \alert{Mô-đun agent} (công cụ chỉ đọc, brief, verifier)
        \end{itemize}
      \end{block}
    \end{column}
    \begin{column}{0.47\textwidth}
      \begin{alertblock}{Đang chờ / chưa làm}
        \begin{itemize}
          \setlength{\itemsep}{1pt}
          \item \alert{Tích hợp agent vào Agentick} (bước tiếp theo)
          \item Registry mô hình, endpoint dự báo
          \item Kiểm tra agent trên dữ liệu Agentick
          \item Ledger/outbox, bảng cảnh báo, dry-run
        \end{itemize}
      \end{alertblock}
      \vspace{0.05cm}
      Agentick là dự án riêng bên ngoài, đã có mã nguồn, dùng để thu dữ liệu
      và chạy thực nghiệm. \alert{Mọi kết quả hôm nay đều từ benchmark TAWOS.}
    \end{column}
  \end{columns}
\end{frame}

\begin{frame}{Roadmap và kế hoạch (tính từ T0 = hôm nay)}
  \scriptsize
  \setlength{\tabcolsep}{3pt}
  \renewcommand{\arraystretch}{1.3}
  \begin{tabular}{L{0.2\textwidth}|L{0.14\textwidth}|L{0.32\textwidth}|L{0.24\textwidth}}
    \textbf{Giai đoạn} & \textbf{Thời gian} & \textbf{Công việc} & \textbf{Tiêu chí xong} \\
    \hline
    1. Tích hợp Agentick & T0 $\to$ +2 tuần & Nối dữ liệu vào pipeline as-of; endpoint dự báo; bảng cảnh báo, ledger, dry-run & Agentick hiển thị issue rủi ro + brief; log có phiên bản \\
    2. Hoàn thiện kiểm tra agent & song song bước 1 & Chạy agent trên issue được chọn từ Agentick; kiểm tra trích dẫn, số liệu & Agent chạy trong luồng cảnh báo, brief lưu kèm trace \\
    3. Thực nghiệm & +2 $\to$ +5 tuần & Shadow mode, hiệu chuẩn lại, so với quy tắc & Recall, precision, lead time trên dữ liệu thật \\
    4. Báo cáo cuối & 2 tuần cuối & Tổng hợp, hoàn thiện bài báo và ứng dụng & Báo cáo, mã, dữ liệu tái lập \\
  \end{tabular}

  \vspace{0.2cm}
  \footnotesize
  \alert{Shadow mode:} dự báo và brief được ghi log nhưng chưa gửi cho người dùng.
  Chỉ sang bước 3 khi dữ liệu Agentick tái dựng được phạm vi cam kết và trạng thái
  hoàn thành, theo cùng tiêu chí đã dùng với TAWOS.
\end{frame}

\begin{frame}{Tóm tắt}
  \small
  \begin{itemize}
    \setlength{\itemsep}{4pt}
    \item Lịch sử thực thi \alert{cải thiện} khả năng phát hiện issue sắp trễ
          (+8,90 / +9,96 điểm tại 50\%), giữ được trên dự án chưa thấy.
    \item Phần lớn lợi thế đến từ việc nhận biết issue đã xong. Trên issue còn
          mở, mức cải thiện \alert{nhỏ nhưng vẫn có ý nghĩa}.
    \item Mô-đun agent đã xây xong: tự lấy bằng chứng tại thời điểm dự báo,
          đề xuất cảnh báo và viết brief \alert{trích dẫn sự kiện có thật}.
    \item \alert{Bước tiếp theo:} tích hợp vào Agentick, kiểm tra agent, thực
          nghiệm shadow mode trên dữ liệu thật.
  \end{itemize}
\end{frame}

\begin{frame}{Tài liệu tham khảo}
  \scriptsize
  \begin{thebibliography}{9}
    \setbeamertemplate{bibliography item}[text]
    \bibitem{tu} F.~Tu, J.~Zhu, Q.~Zheng, and M.~Zhou, ``Be careful of when: An empirical study on time-related misuse of issue tracking data,'' in \emph{Proc. ESEC/FSE}, 2018.
    \bibitem{catboost} L.~Prokhorenkova \emph{et al.}, ``CatBoost: Unbiased boosting with categorical features,'' in \emph{Proc. NeurIPS}, 2018.
    \bibitem{agentdojo} E.~Debenedetti \emph{et al.}, ``AgentDojo: A dynamic environment to evaluate prompt injection attacks and defenses for LLM agents,'' arXiv:2406.13352, 2024.
    \bibitem{tawos} V.~Tawosi, A.~Al-Subaihin, R.~Moussa, and F.~Sarro, ``A versatile dataset of agile open source software projects,'' in \emph{Proc. MSR}, 2022, pp.~707--711.
    \bibitem{verenich} I.~Verenich \emph{et al.}, ``Survey and cross-benchmark comparison of remaining time prediction methods in business process monitoring,'' \emph{ACM TIST}, 2019.
    \bibitem{kubrak} M.~Kubrak \emph{et al.}, ``Prescriptive process monitoring: Quo vadis?,'' \emph{PeerJ Comput. Sci.}, 2022.
  \end{thebibliography}
\end{frame}
